\section{評価実験の方法} \label{sec4}
本実験の目的は，RAG構成を適応的に選択するRMの有効性を評価することである．
実験はQAタスクにおいて，(i) RMの性能評価，(ii) 選択されたRAG構成を用いたRAGの性能評価および，(iii) ユーザを考慮したRAG構成選択の性能評価の三つの観点から構成する．
本章では各評価実験の方法について述べる．

\subsection{評価項目} \label{sec4.0}
評価項目(i)ではRMによる予測順位と正解順位の一致度を検証する．
評価項目(ii)ではRMが選択したRAG構成を用いてRAGを実行し，その最終的な回答を五つの重み設定および各ベースライン手法と比較評価する．
本実験ではRAG構成を上位1個（提案@1）もしくは5個（提案@5）選択し，5個の場合は平均値を算出する．
% 本実験ではRAG構成を上位1個と5個選択し，5個の場合は平均値を算出する．
評価項目(iii)では，提案手法が提示する上位K個のRAG構成を，ユーザが存在するという実運用に近い状況での利用を想定し，その有効性を検証する．
具体的には，ユーザが上位K個の候補それぞれについてRAGを実行し，得られた性能指標に基づいて最終的に一つのRAG構成を採用するという運用フローを想定する．

実験手順として，まず各重視方針（適合率重視，再現率重視，$F_1$値重視，時間重視）に対応する重み設定（第~\ref{sec4.3}節）を用いてRMにより各質問の上位K個の候補構成を抽出する．次に，各質問に対してその上位K個の候補を用いて推論し，それぞれの出力を正解データと比較して適合率，再現率，$F_1$値，および実行時間を算出する．
その後，運用上の意思決定基準を模した「重視方針」に従い，K個の結果の中から当該方針に対応する指標が最良となる一つの候補を選択する．
選択された候補に対応する指標値をその質問の代表値とし，全質問にわたる平均を計算することで，方針別の性能を算出する．
具体的には，「$F_1$値重視」では，$W_{\mathrm{f_1}}=1.0$ で抽出された上位K個のうち実測$F_1$値が最良となるRAG構成を一つ採用する．
同様に，「適合率重視」は $W_{\mathrm{pre}}=1.0$ での上位K個から実測適合率最良構成，「再現率重視」は $W_{\mathrm{rec}}=1.0$ での上位K個から実測再現率最良構成，および「時間重視」は $W_{\mathrm{time}}=1.0$ での上位K個から実測実行時間最小構成をそれぞれ選択する．

\subsection{評価用データセット} \label{sec4.1}
本実験では，SQuAD~\cite{rajpurkar16}，HotpotQA~\cite{yang18}，2WikiMultihopQA~\cite{ho20}，MuSiQue~\cite{trivedi22}の四つのQAデータセットを用いた．
各データセットは，質問，対応する文書，および正解から構成される．
多様な難易度で評価するため，四つのデータセットを統合して用いた．
各データセットから，訓練用として100事例，評価用として20事例をランダムに抽出した．
全データセットの合計は，訓練データが400事例，評価データが80事例である．

\subsection{ランキングモデル} \label{sec4.2}
本実験のランキングタスクにおける学習モデルとして，LightGBM~\cite{KeGuolin17}，XGBoost~\cite{ChenTianqi16}，および CatBoost~\cite{Prokhorenkova18} を採用した．
LightGBM および XGBoost は，非線形な特徴間の相互作用を学習可能であるという特徴を持つ．
CatBoost は，カテゴリ変数のエンコーディングに特化したアルゴリズムである．
なお，評価項目(ii)と評価項目(iii)では，評価項目(i)において最も高い性能を示した学習モデルを選択し，評価に用いた．

\subsection{提案手法の設定} \label{sec4.3}
本研究の実装詳細および評価における重み付けの設定について述べる．
評価項目(ii)と(iii)においては，式~\eqref{eq1}の重み付けが異なる5種類の提案手法を設定した．
第一の設定は精度とコストのバランスを重視し，$W_{\mathrm{f_1}}=W_{\mathrm{time}}=0.5$ とした．
第二の設定は適合率を重視し，$W_{\mathrm{pre}}=1.0$ ，第三の設定は再現率を重視し，$W_{\mathrm{rec}}=1.0$ とした．
第四の設定は $F_1$ 値重視で，$W_{\mathrm{f_1}}=1.0$ を重視し，第五の設定は時間を重視し，$W_{\mathrm{time}}=1.0$ とした．

\subsection{ベースライン手法} \label{sec4.5}
本実験では，比較対象として複数のベースライン手法を設定した．
評価項目(i)では，ランダムに生成したランキングをベースライン手法とし，異なるシード値で10回試行した平均値を算出した．
評価項目(ii)では，以下の五つの設定を定義した．
% 第一に，初期実装で採用されやすい構成として，低コストで構成された``NaiveRAG，stuff，250，50，llama3.1:8b，nomic-embed-text''を全質問に適用する「標準設定」を設けた．
第一に，LangChainのデフォルト設定に基づき開発者が初期段階で採用する構成を「標準設定」とした．
この設定では，すべての質問に対して「NaiveRAG，stuff，250，50，llama3.1:8b，nomic-embed-text」を適用し，ローカル環境における運用の基準点とした．
第二に，提案手法による$F_1$値と実行時間の評価において頻出した構成を「最頻出設定」とした．
具体的には，「NaiveRAG，stuff，500，50，llama3.1:8b，nomic-embed-text」に固定した構成を指す．
第三に，全96通りの候補構成空間から無作為に構成を選択する「ランダム選択」を設定した．これは各質問に対して候補を一様無作為に割り当てた場合の期待値であり，全7,680試行（80質問 $\times$ 96構成）の全体平均として算出した．
第四に，先行研究の代表的な適応的RAG手法であるAdaptive-RAG~\cite{Jeong24}を設定した．Adaptive-RAGは，質問の複雑度分類に基づいてマクロなRAG検索戦略（No-Retrieval / Single-Step / Multi-Step）を切り替えるルーティング手法である．本実験では，公式実装に基づき訓練された質問の3クラス分類器を用い，出力クラスに応じて代表的な構成（No-Retrieval：実行時間最小構成，Single-Step：標準設定，Multi-Step：GraphRAG構成）を割り当ててシミュレーションした．
最後に，候補集合から評価指標が最大となる設定を事後的に選択する「オラクル」を定義し，これを理論上の上限とした．

\subsection{評価指標} \label{sec4.6}
RMの評価指標としてNDCG@K~\cite{JarvelinKalervo02}およびHit@K~\footnote{\url{https://docs.shaped.ai/docs/metrics/hr/}}を用いた．
NDCG@Kは，適切な候補が上位に位置するほど高い値となる指標であり，上位候補の品質を定量的に評価できる．
Hit@Kは，上位K件の候補の中にオラクル構成が含まれる質問の割合であり，実用的なランキング性能を直接的に示す．

RAGの回答性能を評価するため，意味類似度指標であるBERTScore~\cite{zhang20}に基づく適合率，再現率，$F_1$値に加え，QAタスクの事実的・字面的な正解度を評価する標準指標としてExact Match (EM)およびToken-level $F_1$値，ならびに実行時間を採用する．
また，提案手法の優位性の統計的信頼性を保証するため，提案手法と主要なベースライン（特にAdaptive-RAG）の回答精度（$F_1$値）の差について，対応のある有意差検定であるWilcoxon符号付き順位検定を行う．
実行時間は質問と文書が与えられてから回答を出力するまでの時間を計測した．
これらの指標を用い，RMと重みの有効性を検証する．

\subsection{実装環境} \label{sec4.7}
NaiveRAGの実装にはLangChain~\footnote{\url{https://www.langchain.com/}}を，
グラフベースRAGにはGitHub~\footnote{\url{https://github.com/HKUDS/LightRAG/}}~\footnote{\url{https://github.com/microsoft/graphrag/}}で公開されている実装を用いた．
特徴量抽出において，トークン数，依存構造木の深さ，固有表現の抽出にはspaCy~\footnote{\url{https://spacy.io/}}を，キーワードの抽出にはYAKE~\footnote{\url{http://yake.inesctec.pt/}}を用いる．
% 生成モデルと埋め込みモデルについては，Ollama~\footnote{\url{https://ollama.com/}}を用いて実装した．
% GPUにはNVIDIA A100を用い，提案手法で選択しないハイパーパラメータはデフォルト値に固定した．
% また，すべての実験で乱数シードを固定し，再現性を確保した．
GPUにはNVIDIA A100を用い，すべての実験で乱数シードを固定し，再現性を確保した．
なお，RMの学習と推論のためのソースコードは，GitHub~\footnote{\url{https://github.com/ke-lab-it-agu/Adaptive-RAG-Configuration-Selection/}}で公開している． 